\documentclass[12pt,a4paper]{article}
\usepackage[utf8]{inputenc}
\usepackage[indonesian]{babel}
\usepackage{amsmath}
\usepackage{amsfonts}
\usepackage{amssymb}
\usepackage{geometry}

\geometry{margin=2.5cm}

\title{Penyelesaian Volume Benda dengan Metode Shell}
\author{Ahmad Fakhrul Bawani}
\date{}

\begin{document}

\maketitle

\section*{Soal}

Diketahui:
\begin{itemize}
  \item Fungsi pembatas: $x = 4 - 2y^2$
  \item Batas-batas: $y = 0$ dan $y = \sqrt{2}$
  \item Sumbu putar: Sumbu-$x$
\end{itemize}
Tentukan volumenya

\subsection*{1. Rumus Umum}
Karena diputar mengelilingi sumbu-$x$, fungsi diintegrasikan terhadap $y$:
\begin{equation*}
  V = 2\pi \int_{c}^{d} (\text{jari-jari}) \cdot (\text{tinggi}) \, dy
\end{equation*}
Dimana jari-jari tabung adalah $y$ dan tinggi tabung adalah $x = 4 - 2y^2$.

\subsection*{2. Solusi Integrasi}
Substitusikan komponen yang diketahui ke dalam rumus:
\begin{align*}
  V &= 2\pi \int_{0}^{\sqrt{2}} y(4 - 2y^2) \, dy \\
  V &= 2\pi \int_{0}^{\sqrt{2}} (4y - 2y^3) \, dy \\
  V &= 2\pi \left[ 2y^2 - \frac{1}{2}y^4 \right]_{0}^{\sqrt{2}} \\
  V &= 2\pi \left( \left[ 2(\sqrt{2})^2 - \frac{1}{2}(\sqrt{2})^4 \right] - 0 \right) \\
  V &= 2\pi \left( 2(2) - \frac{1}{2}(4) \right) \\
  V &= 2\pi (4 - 2) \\
  V &= 4\pi
\end{align*}

Jadi, volume benda putar tersebut adalah \textbf{$4\pi$ satuan volume}.

\end{document}